\documentclass[11pt]{article}

    \usepackage[breakable]{tcolorbox}
    \usepackage{parskip} % Stop auto-indenting (to mimic markdown behaviour)
    

    % Basic figure setup, for now with no caption control since it's done
    % automatically by Pandoc (which extracts ![](path) syntax from Markdown).
    \usepackage{graphicx}
    % Keep aspect ratio if custom image width or height is specified
    \setkeys{Gin}{keepaspectratio}
    % Maintain compatibility with old templates. Remove in nbconvert 6.0
    \let\Oldincludegraphics\includegraphics
    % Ensure that by default, figures have no caption (until we provide a
    % proper Figure object with a Caption API and a way to capture that
    % in the conversion process - todo).
    \usepackage{caption}
    \DeclareCaptionFormat{nocaption}{}
    \captionsetup{format=nocaption,aboveskip=0pt,belowskip=0pt}

    \usepackage{float}
    \floatplacement{figure}{H} % forces figures to be placed at the correct location
    \usepackage{xcolor} % Allow colors to be defined
    \usepackage{enumerate} % Needed for markdown enumerations to work
    \usepackage{geometry} % Used to adjust the document margins
    \usepackage{amsmath} % Equations
    \usepackage{amssymb} % Equations
    \usepackage{textcomp} % defines textquotesingle
    % Hack from http://tex.stackexchange.com/a/47451/13684:
    \AtBeginDocument{%
        \def\PYZsq{\textquotesingle}% Upright quotes in Pygmentized code
    }
    \usepackage{upquote} % Upright quotes for verbatim code
    \usepackage{eurosym} % defines \euro

    \usepackage{iftex}
    \ifPDFTeX
        \usepackage[T1]{fontenc}
        \IfFileExists{alphabeta.sty}{
              \usepackage{alphabeta}
          }{
              \usepackage[mathletters]{ucs}
              \usepackage[utf8x]{inputenc}
          }
    \else
        \usepackage{fontspec}
        \usepackage{unicode-math}
    \fi

    \usepackage{fancyvrb} % verbatim replacement that allows latex
    \usepackage{grffile} % extends the file name processing of package graphics
                         % to support a larger range
    \makeatletter % fix for old versions of grffile with XeLaTeX
    \@ifpackagelater{grffile}{2019/11/01}
    {
      % Do nothing on new versions
    }
    {
      \def\Gread@@xetex#1{%
        \IfFileExists{"\Gin@base".bb}%
        {\Gread@eps{\Gin@base.bb}}%
        {\Gread@@xetex@aux#1}%
      }
    }
    \makeatother
    \usepackage[Export]{adjustbox} % Used to constrain images to a maximum size
    \adjustboxset{max size={0.9\linewidth}{0.9\paperheight}}

    % The hyperref package gives us a pdf with properly built
    % internal navigation ('pdf bookmarks' for the table of contents,
    % internal cross-reference links, web links for URLs, etc.)
    \usepackage{hyperref}
    % The default LaTeX title has an obnoxious amount of whitespace. By default,
    % titling removes some of it. It also provides customization options.
    \usepackage{titling}
    \usepackage{longtable} % longtable support required by pandoc >1.10
    \usepackage{booktabs}  % table support for pandoc > 1.12.2
    \usepackage{array}     % table support for pandoc >= 2.11.3
    \usepackage{calc}      % table minipage width calculation for pandoc >= 2.11.1
    \usepackage[inline]{enumitem} % IRkernel/repr support (it uses the enumerate* environment)
    \usepackage[normalem]{ulem} % ulem is needed to support strikethroughs (\sout)
                                % normalem makes italics be italics, not underlines
    \usepackage{soul}      % strikethrough (\st) support for pandoc >= 3.0.0
    \usepackage{mathrsfs}
    

    
    % Colors for the hyperref package
    \definecolor{urlcolor}{rgb}{0,.145,.698}
    \definecolor{linkcolor}{rgb}{.71,0.21,0.01}
    \definecolor{citecolor}{rgb}{.12,.54,.11}

    % ANSI colors
    \definecolor{ansi-black}{HTML}{3E424D}
    \definecolor{ansi-black-intense}{HTML}{282C36}
    \definecolor{ansi-red}{HTML}{E75C58}
    \definecolor{ansi-red-intense}{HTML}{B22B31}
    \definecolor{ansi-green}{HTML}{00A250}
    \definecolor{ansi-green-intense}{HTML}{007427}
    \definecolor{ansi-yellow}{HTML}{DDB62B}
    \definecolor{ansi-yellow-intense}{HTML}{B27D12}
    \definecolor{ansi-blue}{HTML}{208FFB}
    \definecolor{ansi-blue-intense}{HTML}{0065CA}
    \definecolor{ansi-magenta}{HTML}{D160C4}
    \definecolor{ansi-magenta-intense}{HTML}{A03196}
    \definecolor{ansi-cyan}{HTML}{60C6C8}
    \definecolor{ansi-cyan-intense}{HTML}{258F8F}
    \definecolor{ansi-white}{HTML}{C5C1B4}
    \definecolor{ansi-white-intense}{HTML}{A1A6B2}
    \definecolor{ansi-default-inverse-fg}{HTML}{FFFFFF}
    \definecolor{ansi-default-inverse-bg}{HTML}{000000}

    % common color for the border for error outputs.
    \definecolor{outerrorbackground}{HTML}{FFDFDF}

    % commands and environments needed by pandoc snippets
    % extracted from the output of `pandoc -s`
    \providecommand{\tightlist}{%
      \setlength{\itemsep}{0pt}\setlength{\parskip}{0pt}}
    \DefineVerbatimEnvironment{Highlighting}{Verbatim}{commandchars=\\\{\}}
    % Add ',fontsize=\small' for more characters per line
    \newenvironment{Shaded}{}{}
    \newcommand{\KeywordTok}[1]{\textcolor[rgb]{0.00,0.44,0.13}{\textbf{{#1}}}}
    \newcommand{\DataTypeTok}[1]{\textcolor[rgb]{0.56,0.13,0.00}{{#1}}}
    \newcommand{\DecValTok}[1]{\textcolor[rgb]{0.25,0.63,0.44}{{#1}}}
    \newcommand{\BaseNTok}[1]{\textcolor[rgb]{0.25,0.63,0.44}{{#1}}}
    \newcommand{\FloatTok}[1]{\textcolor[rgb]{0.25,0.63,0.44}{{#1}}}
    \newcommand{\CharTok}[1]{\textcolor[rgb]{0.25,0.44,0.63}{{#1}}}
    \newcommand{\StringTok}[1]{\textcolor[rgb]{0.25,0.44,0.63}{{#1}}}
    \newcommand{\CommentTok}[1]{\textcolor[rgb]{0.38,0.63,0.69}{\textit{{#1}}}}
    \newcommand{\OtherTok}[1]{\textcolor[rgb]{0.00,0.44,0.13}{{#1}}}
    \newcommand{\AlertTok}[1]{\textcolor[rgb]{1.00,0.00,0.00}{\textbf{{#1}}}}
    \newcommand{\FunctionTok}[1]{\textcolor[rgb]{0.02,0.16,0.49}{{#1}}}
    \newcommand{\RegionMarkerTok}[1]{{#1}}
    \newcommand{\ErrorTok}[1]{\textcolor[rgb]{1.00,0.00,0.00}{\textbf{{#1}}}}
    \newcommand{\NormalTok}[1]{{#1}}

    % Additional commands for more recent versions of Pandoc
    \newcommand{\ConstantTok}[1]{\textcolor[rgb]{0.53,0.00,0.00}{{#1}}}
    \newcommand{\SpecialCharTok}[1]{\textcolor[rgb]{0.25,0.44,0.63}{{#1}}}
    \newcommand{\VerbatimStringTok}[1]{\textcolor[rgb]{0.25,0.44,0.63}{{#1}}}
    \newcommand{\SpecialStringTok}[1]{\textcolor[rgb]{0.73,0.40,0.53}{{#1}}}
    \newcommand{\ImportTok}[1]{{#1}}
    \newcommand{\DocumentationTok}[1]{\textcolor[rgb]{0.73,0.13,0.13}{\textit{{#1}}}}
    \newcommand{\AnnotationTok}[1]{\textcolor[rgb]{0.38,0.63,0.69}{\textbf{\textit{{#1}}}}}
    \newcommand{\CommentVarTok}[1]{\textcolor[rgb]{0.38,0.63,0.69}{\textbf{\textit{{#1}}}}}
    \newcommand{\VariableTok}[1]{\textcolor[rgb]{0.10,0.09,0.49}{{#1}}}
    \newcommand{\ControlFlowTok}[1]{\textcolor[rgb]{0.00,0.44,0.13}{\textbf{{#1}}}}
    \newcommand{\OperatorTok}[1]{\textcolor[rgb]{0.40,0.40,0.40}{{#1}}}
    \newcommand{\BuiltInTok}[1]{{#1}}
    \newcommand{\ExtensionTok}[1]{{#1}}
    \newcommand{\PreprocessorTok}[1]{\textcolor[rgb]{0.74,0.48,0.00}{{#1}}}
    \newcommand{\AttributeTok}[1]{\textcolor[rgb]{0.49,0.56,0.16}{{#1}}}
    \newcommand{\InformationTok}[1]{\textcolor[rgb]{0.38,0.63,0.69}{\textbf{\textit{{#1}}}}}
    \newcommand{\WarningTok}[1]{\textcolor[rgb]{0.38,0.63,0.69}{\textbf{\textit{{#1}}}}}
    \makeatletter
    \newsavebox\pandoc@box
    \newcommand*\pandocbounded[1]{%
      \sbox\pandoc@box{#1}%
      % scaling factors for width and height
      \Gscale@div\@tempa\textheight{\dimexpr\ht\pandoc@box+\dp\pandoc@box\relax}%
      \Gscale@div\@tempb\linewidth{\wd\pandoc@box}%
      % select the smaller of both
      \ifdim\@tempb\p@<\@tempa\p@
        \let\@tempa\@tempb
      \fi
      % scaling accordingly (\@tempa < 1)
      \ifdim\@tempa\p@<\p@
        \scalebox{\@tempa}{\usebox\pandoc@box}%
      % scaling not needed, use as it is
      \else
        \usebox{\pandoc@box}%
      \fi
    }
    \makeatother

    % Define a nice break command that doesn't care if a line doesn't already
    % exist.
    \def\br{\hspace*{\fill} \\* }
    % Math Jax compatibility definitions
    \def\gt{>}
    \def\lt{<}
    \let\Oldtex\TeX
    \let\Oldlatex\LaTeX
    \renewcommand{\TeX}{\textrm{\Oldtex}}
    \renewcommand{\LaTeX}{\textrm{\Oldlatex}}
    % Document parameters
    % Document title
    \title{Ch05-resample-lab}
    
    
    
    
    
    
    
% Pygments definitions
\makeatletter
\def\PY@reset{\let\PY@it=\relax \let\PY@bf=\relax%
    \let\PY@ul=\relax \let\PY@tc=\relax%
    \let\PY@bc=\relax \let\PY@ff=\relax}
\def\PY@tok#1{\csname PY@tok@#1\endcsname}
\def\PY@toks#1+{\ifx\relax#1\empty\else%
    \PY@tok{#1}\expandafter\PY@toks\fi}
\def\PY@do#1{\PY@bc{\PY@tc{\PY@ul{%
    \PY@it{\PY@bf{\PY@ff{#1}}}}}}}
\def\PY#1#2{\PY@reset\PY@toks#1+\relax+\PY@do{#2}}

\@namedef{PY@tok@w}{\def\PY@tc##1{\textcolor[rgb]{0.73,0.73,0.73}{##1}}}
\@namedef{PY@tok@c}{\let\PY@it=\textit\def\PY@tc##1{\textcolor[rgb]{0.24,0.48,0.48}{##1}}}
\@namedef{PY@tok@cp}{\def\PY@tc##1{\textcolor[rgb]{0.61,0.40,0.00}{##1}}}
\@namedef{PY@tok@k}{\let\PY@bf=\textbf\def\PY@tc##1{\textcolor[rgb]{0.00,0.50,0.00}{##1}}}
\@namedef{PY@tok@kp}{\def\PY@tc##1{\textcolor[rgb]{0.00,0.50,0.00}{##1}}}
\@namedef{PY@tok@kt}{\def\PY@tc##1{\textcolor[rgb]{0.69,0.00,0.25}{##1}}}
\@namedef{PY@tok@o}{\def\PY@tc##1{\textcolor[rgb]{0.40,0.40,0.40}{##1}}}
\@namedef{PY@tok@ow}{\let\PY@bf=\textbf\def\PY@tc##1{\textcolor[rgb]{0.67,0.13,1.00}{##1}}}
\@namedef{PY@tok@nb}{\def\PY@tc##1{\textcolor[rgb]{0.00,0.50,0.00}{##1}}}
\@namedef{PY@tok@nf}{\def\PY@tc##1{\textcolor[rgb]{0.00,0.00,1.00}{##1}}}
\@namedef{PY@tok@nc}{\let\PY@bf=\textbf\def\PY@tc##1{\textcolor[rgb]{0.00,0.00,1.00}{##1}}}
\@namedef{PY@tok@nn}{\let\PY@bf=\textbf\def\PY@tc##1{\textcolor[rgb]{0.00,0.00,1.00}{##1}}}
\@namedef{PY@tok@ne}{\let\PY@bf=\textbf\def\PY@tc##1{\textcolor[rgb]{0.80,0.25,0.22}{##1}}}
\@namedef{PY@tok@nv}{\def\PY@tc##1{\textcolor[rgb]{0.10,0.09,0.49}{##1}}}
\@namedef{PY@tok@no}{\def\PY@tc##1{\textcolor[rgb]{0.53,0.00,0.00}{##1}}}
\@namedef{PY@tok@nl}{\def\PY@tc##1{\textcolor[rgb]{0.46,0.46,0.00}{##1}}}
\@namedef{PY@tok@ni}{\let\PY@bf=\textbf\def\PY@tc##1{\textcolor[rgb]{0.44,0.44,0.44}{##1}}}
\@namedef{PY@tok@na}{\def\PY@tc##1{\textcolor[rgb]{0.41,0.47,0.13}{##1}}}
\@namedef{PY@tok@nt}{\let\PY@bf=\textbf\def\PY@tc##1{\textcolor[rgb]{0.00,0.50,0.00}{##1}}}
\@namedef{PY@tok@nd}{\def\PY@tc##1{\textcolor[rgb]{0.67,0.13,1.00}{##1}}}
\@namedef{PY@tok@s}{\def\PY@tc##1{\textcolor[rgb]{0.73,0.13,0.13}{##1}}}
\@namedef{PY@tok@sd}{\let\PY@it=\textit\def\PY@tc##1{\textcolor[rgb]{0.73,0.13,0.13}{##1}}}
\@namedef{PY@tok@si}{\let\PY@bf=\textbf\def\PY@tc##1{\textcolor[rgb]{0.64,0.35,0.47}{##1}}}
\@namedef{PY@tok@se}{\let\PY@bf=\textbf\def\PY@tc##1{\textcolor[rgb]{0.67,0.36,0.12}{##1}}}
\@namedef{PY@tok@sr}{\def\PY@tc##1{\textcolor[rgb]{0.64,0.35,0.47}{##1}}}
\@namedef{PY@tok@ss}{\def\PY@tc##1{\textcolor[rgb]{0.10,0.09,0.49}{##1}}}
\@namedef{PY@tok@sx}{\def\PY@tc##1{\textcolor[rgb]{0.00,0.50,0.00}{##1}}}
\@namedef{PY@tok@m}{\def\PY@tc##1{\textcolor[rgb]{0.40,0.40,0.40}{##1}}}
\@namedef{PY@tok@gh}{\let\PY@bf=\textbf\def\PY@tc##1{\textcolor[rgb]{0.00,0.00,0.50}{##1}}}
\@namedef{PY@tok@gu}{\let\PY@bf=\textbf\def\PY@tc##1{\textcolor[rgb]{0.50,0.00,0.50}{##1}}}
\@namedef{PY@tok@gd}{\def\PY@tc##1{\textcolor[rgb]{0.63,0.00,0.00}{##1}}}
\@namedef{PY@tok@gi}{\def\PY@tc##1{\textcolor[rgb]{0.00,0.52,0.00}{##1}}}
\@namedef{PY@tok@gr}{\def\PY@tc##1{\textcolor[rgb]{0.89,0.00,0.00}{##1}}}
\@namedef{PY@tok@ge}{\let\PY@it=\textit}
\@namedef{PY@tok@gs}{\let\PY@bf=\textbf}
\@namedef{PY@tok@ges}{\let\PY@bf=\textbf\let\PY@it=\textit}
\@namedef{PY@tok@gp}{\let\PY@bf=\textbf\def\PY@tc##1{\textcolor[rgb]{0.00,0.00,0.50}{##1}}}
\@namedef{PY@tok@go}{\def\PY@tc##1{\textcolor[rgb]{0.44,0.44,0.44}{##1}}}
\@namedef{PY@tok@gt}{\def\PY@tc##1{\textcolor[rgb]{0.00,0.27,0.87}{##1}}}
\@namedef{PY@tok@err}{\def\PY@bc##1{{\setlength{\fboxsep}{\string -\fboxrule}\fcolorbox[rgb]{1.00,0.00,0.00}{1,1,1}{\strut ##1}}}}
\@namedef{PY@tok@kc}{\let\PY@bf=\textbf\def\PY@tc##1{\textcolor[rgb]{0.00,0.50,0.00}{##1}}}
\@namedef{PY@tok@kd}{\let\PY@bf=\textbf\def\PY@tc##1{\textcolor[rgb]{0.00,0.50,0.00}{##1}}}
\@namedef{PY@tok@kn}{\let\PY@bf=\textbf\def\PY@tc##1{\textcolor[rgb]{0.00,0.50,0.00}{##1}}}
\@namedef{PY@tok@kr}{\let\PY@bf=\textbf\def\PY@tc##1{\textcolor[rgb]{0.00,0.50,0.00}{##1}}}
\@namedef{PY@tok@bp}{\def\PY@tc##1{\textcolor[rgb]{0.00,0.50,0.00}{##1}}}
\@namedef{PY@tok@fm}{\def\PY@tc##1{\textcolor[rgb]{0.00,0.00,1.00}{##1}}}
\@namedef{PY@tok@vc}{\def\PY@tc##1{\textcolor[rgb]{0.10,0.09,0.49}{##1}}}
\@namedef{PY@tok@vg}{\def\PY@tc##1{\textcolor[rgb]{0.10,0.09,0.49}{##1}}}
\@namedef{PY@tok@vi}{\def\PY@tc##1{\textcolor[rgb]{0.10,0.09,0.49}{##1}}}
\@namedef{PY@tok@vm}{\def\PY@tc##1{\textcolor[rgb]{0.10,0.09,0.49}{##1}}}
\@namedef{PY@tok@sa}{\def\PY@tc##1{\textcolor[rgb]{0.73,0.13,0.13}{##1}}}
\@namedef{PY@tok@sb}{\def\PY@tc##1{\textcolor[rgb]{0.73,0.13,0.13}{##1}}}
\@namedef{PY@tok@sc}{\def\PY@tc##1{\textcolor[rgb]{0.73,0.13,0.13}{##1}}}
\@namedef{PY@tok@dl}{\def\PY@tc##1{\textcolor[rgb]{0.73,0.13,0.13}{##1}}}
\@namedef{PY@tok@s2}{\def\PY@tc##1{\textcolor[rgb]{0.73,0.13,0.13}{##1}}}
\@namedef{PY@tok@sh}{\def\PY@tc##1{\textcolor[rgb]{0.73,0.13,0.13}{##1}}}
\@namedef{PY@tok@s1}{\def\PY@tc##1{\textcolor[rgb]{0.73,0.13,0.13}{##1}}}
\@namedef{PY@tok@mb}{\def\PY@tc##1{\textcolor[rgb]{0.40,0.40,0.40}{##1}}}
\@namedef{PY@tok@mf}{\def\PY@tc##1{\textcolor[rgb]{0.40,0.40,0.40}{##1}}}
\@namedef{PY@tok@mh}{\def\PY@tc##1{\textcolor[rgb]{0.40,0.40,0.40}{##1}}}
\@namedef{PY@tok@mi}{\def\PY@tc##1{\textcolor[rgb]{0.40,0.40,0.40}{##1}}}
\@namedef{PY@tok@il}{\def\PY@tc##1{\textcolor[rgb]{0.40,0.40,0.40}{##1}}}
\@namedef{PY@tok@mo}{\def\PY@tc##1{\textcolor[rgb]{0.40,0.40,0.40}{##1}}}
\@namedef{PY@tok@ch}{\let\PY@it=\textit\def\PY@tc##1{\textcolor[rgb]{0.24,0.48,0.48}{##1}}}
\@namedef{PY@tok@cm}{\let\PY@it=\textit\def\PY@tc##1{\textcolor[rgb]{0.24,0.48,0.48}{##1}}}
\@namedef{PY@tok@cpf}{\let\PY@it=\textit\def\PY@tc##1{\textcolor[rgb]{0.24,0.48,0.48}{##1}}}
\@namedef{PY@tok@c1}{\let\PY@it=\textit\def\PY@tc##1{\textcolor[rgb]{0.24,0.48,0.48}{##1}}}
\@namedef{PY@tok@cs}{\let\PY@it=\textit\def\PY@tc##1{\textcolor[rgb]{0.24,0.48,0.48}{##1}}}

\def\PYZbs{\char`\\}
\def\PYZus{\char`\_}
\def\PYZob{\char`\{}
\def\PYZcb{\char`\}}
\def\PYZca{\char`\^}
\def\PYZam{\char`\&}
\def\PYZlt{\char`\<}
\def\PYZgt{\char`\>}
\def\PYZsh{\char`\#}
\def\PYZpc{\char`\%}
\def\PYZdl{\char`\$}
\def\PYZhy{\char`\-}
\def\PYZsq{\char`\'}
\def\PYZdq{\char`\"}
\def\PYZti{\char`\~}
% for compatibility with earlier versions
\def\PYZat{@}
\def\PYZlb{[}
\def\PYZrb{]}
\makeatother


    % For linebreaks inside Verbatim environment from package fancyvrb.
    \makeatletter
        \newbox\Wrappedcontinuationbox
        \newbox\Wrappedvisiblespacebox
        \newcommand*\Wrappedvisiblespace {\textcolor{red}{\textvisiblespace}}
        \newcommand*\Wrappedcontinuationsymbol {\textcolor{red}{\llap{\tiny$\m@th\hookrightarrow$}}}
        \newcommand*\Wrappedcontinuationindent {3ex }
        \newcommand*\Wrappedafterbreak {\kern\Wrappedcontinuationindent\copy\Wrappedcontinuationbox}
        % Take advantage of the already applied Pygments mark-up to insert
        % potential linebreaks for TeX processing.
        %        {, <, #, %, $, ' and ": go to next line.
        %        _, }, ^, &, >, - and ~: stay at end of broken line.
        % Use of \textquotesingle for straight quote.
        \newcommand*\Wrappedbreaksatspecials {%
            \def\PYGZus{\discretionary{\char`\_}{\Wrappedafterbreak}{\char`\_}}%
            \def\PYGZob{\discretionary{}{\Wrappedafterbreak\char`\{}{\char`\{}}%
            \def\PYGZcb{\discretionary{\char`\}}{\Wrappedafterbreak}{\char`\}}}%
            \def\PYGZca{\discretionary{\char`\^}{\Wrappedafterbreak}{\char`\^}}%
            \def\PYGZam{\discretionary{\char`\&}{\Wrappedafterbreak}{\char`\&}}%
            \def\PYGZlt{\discretionary{}{\Wrappedafterbreak\char`\<}{\char`\<}}%
            \def\PYGZgt{\discretionary{\char`\>}{\Wrappedafterbreak}{\char`\>}}%
            \def\PYGZsh{\discretionary{}{\Wrappedafterbreak\char`\#}{\char`\#}}%
            \def\PYGZpc{\discretionary{}{\Wrappedafterbreak\char`\%}{\char`\%}}%
            \def\PYGZdl{\discretionary{}{\Wrappedafterbreak\char`\$}{\char`\$}}%
            \def\PYGZhy{\discretionary{\char`\-}{\Wrappedafterbreak}{\char`\-}}%
            \def\PYGZsq{\discretionary{}{\Wrappedafterbreak\textquotesingle}{\textquotesingle}}%
            \def\PYGZdq{\discretionary{}{\Wrappedafterbreak\char`\"}{\char`\"}}%
            \def\PYGZti{\discretionary{\char`\~}{\Wrappedafterbreak}{\char`\~}}%
        }
        % Some characters . , ; ? ! / are not pygmentized.
        % This macro makes them "active" and they will insert potential linebreaks
        \newcommand*\Wrappedbreaksatpunct {%
            \lccode`\~`\.\lowercase{\def~}{\discretionary{\hbox{\char`\.}}{\Wrappedafterbreak}{\hbox{\char`\.}}}%
            \lccode`\~`\,\lowercase{\def~}{\discretionary{\hbox{\char`\,}}{\Wrappedafterbreak}{\hbox{\char`\,}}}%
            \lccode`\~`\;\lowercase{\def~}{\discretionary{\hbox{\char`\;}}{\Wrappedafterbreak}{\hbox{\char`\;}}}%
            \lccode`\~`\:\lowercase{\def~}{\discretionary{\hbox{\char`\:}}{\Wrappedafterbreak}{\hbox{\char`\:}}}%
            \lccode`\~`\?\lowercase{\def~}{\discretionary{\hbox{\char`\?}}{\Wrappedafterbreak}{\hbox{\char`\?}}}%
            \lccode`\~`\!\lowercase{\def~}{\discretionary{\hbox{\char`\!}}{\Wrappedafterbreak}{\hbox{\char`\!}}}%
            \lccode`\~`\/\lowercase{\def~}{\discretionary{\hbox{\char`\/}}{\Wrappedafterbreak}{\hbox{\char`\/}}}%
            \catcode`\.\active
            \catcode`\,\active
            \catcode`\;\active
            \catcode`\:\active
            \catcode`\?\active
            \catcode`\!\active
            \catcode`\/\active
            \lccode`\~`\~
        }
    \makeatother

    \let\OriginalVerbatim=\Verbatim
    \makeatletter
    \renewcommand{\Verbatim}[1][1]{%
        %\parskip\z@skip
        \sbox\Wrappedcontinuationbox {\Wrappedcontinuationsymbol}%
        \sbox\Wrappedvisiblespacebox {\FV@SetupFont\Wrappedvisiblespace}%
        \def\FancyVerbFormatLine ##1{\hsize\linewidth
            \vtop{\raggedright\hyphenpenalty\z@\exhyphenpenalty\z@
                \doublehyphendemerits\z@\finalhyphendemerits\z@
                \strut ##1\strut}%
        }%
        % If the linebreak is at a space, the latter will be displayed as visible
        % space at end of first line, and a continuation symbol starts next line.
        % Stretch/shrink are however usually zero for typewriter font.
        \def\FV@Space {%
            \nobreak\hskip\z@ plus\fontdimen3\font minus\fontdimen4\font
            \discretionary{\copy\Wrappedvisiblespacebox}{\Wrappedafterbreak}
            {\kern\fontdimen2\font}%
        }%

        % Allow breaks at special characters using \PYG... macros.
        \Wrappedbreaksatspecials
        % Breaks at punctuation characters . , ; ? ! and / need catcode=\active
        \OriginalVerbatim[#1,codes*=\Wrappedbreaksatpunct]%
    }
    \makeatother

    % Exact colors from NB
    \definecolor{incolor}{HTML}{303F9F}
    \definecolor{outcolor}{HTML}{D84315}
    \definecolor{cellborder}{HTML}{CFCFCF}
    \definecolor{cellbackground}{HTML}{F7F7F7}

    % prompt
    \makeatletter
    \newcommand{\boxspacing}{\kern\kvtcb@left@rule\kern\kvtcb@boxsep}
    \makeatother
    \newcommand{\prompt}[4]{
        {\ttfamily\llap{{\color{#2}[#3]:\hspace{3pt}#4}}\vspace{-\baselineskip}}
    }
    

    
    % Prevent overflowing lines due to hard-to-break entities
    \sloppy
    % Setup hyperref package
    \hypersetup{
      breaklinks=true,  % so long urls are correctly broken across lines
      colorlinks=true,
      urlcolor=urlcolor,
      linkcolor=linkcolor,
      citecolor=citecolor,
      }
    % Slightly bigger margins than the latex defaults
    
    \geometry{verbose,tmargin=1in,bmargin=1in,lmargin=1in,rmargin=1in}
    
    

\begin{document}
    
    \maketitle
    
    

    
    \hypertarget{cross-validation-and-the-bootstrap}{%
\section{Cross-Validation and the
Bootstrap}\label{cross-validation-and-the-bootstrap}}

\href{https://mybinder.org/v2/gh/intro-stat-learning/ISLP_labs/v2.2?labpath=Ch05-resample-lab.ipynb}{\includesvg{https://mybinder.org/badge_logo.svg}}

    In this lab, we explore the resampling techniques covered in this
chapter. Some of the commands in this lab may take a while to run on
your computer.

We again begin by placing most of our imports at this top level.

    \begin{tcolorbox}[breakable, size=fbox, boxrule=1pt, pad at break*=1mm,colback=cellbackground, colframe=cellborder]
\prompt{In}{incolor}{1}{\boxspacing}
\begin{Verbatim}[commandchars=\\\{\}]
\PY{k+kn}{import}\PY{+w}{ }\PY{n+nn}{numpy}\PY{+w}{ }\PY{k}{as}\PY{+w}{ }\PY{n+nn}{np}
\PY{k+kn}{import}\PY{+w}{ }\PY{n+nn}{statsmodels}\PY{n+nn}{.}\PY{n+nn}{api}\PY{+w}{ }\PY{k}{as}\PY{+w}{ }\PY{n+nn}{sm}
\PY{k+kn}{from}\PY{+w}{ }\PY{n+nn}{ISLP}\PY{+w}{ }\PY{k+kn}{import} \PY{n}{load\PYZus{}data}
\PY{k+kn}{from}\PY{+w}{ }\PY{n+nn}{ISLP}\PY{n+nn}{.}\PY{n+nn}{models}\PY{+w}{ }\PY{k+kn}{import} \PY{p}{(}\PY{n}{ModelSpec} \PY{k}{as} \PY{n}{MS}\PY{p}{,}
                         \PY{n}{summarize}\PY{p}{,}
                         \PY{n}{poly}\PY{p}{)}
\PY{k+kn}{from}\PY{+w}{ }\PY{n+nn}{sklearn}\PY{n+nn}{.}\PY{n+nn}{model\PYZus{}selection}\PY{+w}{ }\PY{k+kn}{import} \PY{n}{train\PYZus{}test\PYZus{}split}
\end{Verbatim}
\end{tcolorbox}

    There are several new imports needed for this lab.

    \begin{tcolorbox}[breakable, size=fbox, boxrule=1pt, pad at break*=1mm,colback=cellbackground, colframe=cellborder]
\prompt{In}{incolor}{2}{\boxspacing}
\begin{Verbatim}[commandchars=\\\{\}]
\PY{k+kn}{from}\PY{+w}{ }\PY{n+nn}{functools}\PY{+w}{ }\PY{k+kn}{import} \PY{n}{partial}
\PY{k+kn}{from}\PY{+w}{ }\PY{n+nn}{sklearn}\PY{n+nn}{.}\PY{n+nn}{model\PYZus{}selection}\PY{+w}{ }\PY{k+kn}{import} \PYZbs{}
     \PY{p}{(}\PY{n}{cross\PYZus{}validate}\PY{p}{,}
      \PY{n}{KFold}\PY{p}{,}
      \PY{n}{ShuffleSplit}\PY{p}{)}
\PY{k+kn}{from}\PY{+w}{ }\PY{n+nn}{sklearn}\PY{n+nn}{.}\PY{n+nn}{base}\PY{+w}{ }\PY{k+kn}{import} \PY{n}{clone}
\PY{k+kn}{from}\PY{+w}{ }\PY{n+nn}{ISLP}\PY{n+nn}{.}\PY{n+nn}{models}\PY{+w}{ }\PY{k+kn}{import} \PY{n}{sklearn\PYZus{}sm}
\end{Verbatim}
\end{tcolorbox}

    \hypertarget{the-validation-set-approach}{%
\subsection{The Validation Set
Approach}\label{the-validation-set-approach}}

We explore the use of the validation set approach in order to estimate
the test error rates that result from fitting various linear models on
the \texttt{Auto} data set.

We use the function \texttt{train\_test\_split()} to split the data into
training and validation sets. As there are 392 observations, we split
into two equal sets of size 196 using the argument
\texttt{test\_size=196}. It is generally a good idea to set a random
seed when performing operations like this that contain an element of
randomness, so that the results obtained can be reproduced precisely at
a later time. We set the random seed of the splitter with the argument
\texttt{random\_state=0}.

    \begin{tcolorbox}[breakable, size=fbox, boxrule=1pt, pad at break*=1mm,colback=cellbackground, colframe=cellborder]
\prompt{In}{incolor}{3}{\boxspacing}
\begin{Verbatim}[commandchars=\\\{\}]
\PY{n}{Auto} \PY{o}{=} \PY{n}{load\PYZus{}data}\PY{p}{(}\PY{l+s+s1}{\PYZsq{}}\PY{l+s+s1}{Auto}\PY{l+s+s1}{\PYZsq{}}\PY{p}{)}
\PY{n}{Auto\PYZus{}train}\PY{p}{,} \PY{n}{Auto\PYZus{}valid} \PY{o}{=} \PY{n}{train\PYZus{}test\PYZus{}split}\PY{p}{(}\PY{n}{Auto}\PY{p}{,}
                                         \PY{n}{test\PYZus{}size}\PY{o}{=}\PY{l+m+mi}{196}\PY{p}{,}
                                         \PY{n}{random\PYZus{}state}\PY{o}{=}\PY{l+m+mi}{0}\PY{p}{)}
\end{Verbatim}
\end{tcolorbox}

    Now we can fit a linear regression using only the observations
corresponding to the training set \texttt{Auto\_train}.

    \begin{tcolorbox}[breakable, size=fbox, boxrule=1pt, pad at break*=1mm,colback=cellbackground, colframe=cellborder]
\prompt{In}{incolor}{4}{\boxspacing}
\begin{Verbatim}[commandchars=\\\{\}]
\PY{n}{hp\PYZus{}mm} \PY{o}{=} \PY{n}{MS}\PY{p}{(}\PY{p}{[}\PY{l+s+s1}{\PYZsq{}}\PY{l+s+s1}{horsepower}\PY{l+s+s1}{\PYZsq{}}\PY{p}{]}\PY{p}{)}
\PY{n}{X\PYZus{}train} \PY{o}{=} \PY{n}{hp\PYZus{}mm}\PY{o}{.}\PY{n}{fit\PYZus{}transform}\PY{p}{(}\PY{n}{Auto\PYZus{}train}\PY{p}{)}
\PY{n}{y\PYZus{}train} \PY{o}{=} \PY{n}{Auto\PYZus{}train}\PY{p}{[}\PY{l+s+s1}{\PYZsq{}}\PY{l+s+s1}{mpg}\PY{l+s+s1}{\PYZsq{}}\PY{p}{]}
\PY{n}{model} \PY{o}{=} \PY{n}{sm}\PY{o}{.}\PY{n}{OLS}\PY{p}{(}\PY{n}{y\PYZus{}train}\PY{p}{,} \PY{n}{X\PYZus{}train}\PY{p}{)}
\PY{n}{results} \PY{o}{=} \PY{n}{model}\PY{o}{.}\PY{n}{fit}\PY{p}{(}\PY{p}{)}
\end{Verbatim}
\end{tcolorbox}

    We now use the \texttt{predict()} method of \texttt{results} evaluated
on the model matrix for this model created using the validation data
set. We also calculate the validation MSE of our model.

    \begin{tcolorbox}[breakable, size=fbox, boxrule=1pt, pad at break*=1mm,colback=cellbackground, colframe=cellborder]
\prompt{In}{incolor}{5}{\boxspacing}
\begin{Verbatim}[commandchars=\\\{\}]
\PY{n}{X\PYZus{}valid} \PY{o}{=} \PY{n}{hp\PYZus{}mm}\PY{o}{.}\PY{n}{transform}\PY{p}{(}\PY{n}{Auto\PYZus{}valid}\PY{p}{)}
\PY{n}{y\PYZus{}valid} \PY{o}{=} \PY{n}{Auto\PYZus{}valid}\PY{p}{[}\PY{l+s+s1}{\PYZsq{}}\PY{l+s+s1}{mpg}\PY{l+s+s1}{\PYZsq{}}\PY{p}{]}
\PY{n}{valid\PYZus{}pred} \PY{o}{=} \PY{n}{results}\PY{o}{.}\PY{n}{predict}\PY{p}{(}\PY{n}{X\PYZus{}valid}\PY{p}{)}
\PY{n}{np}\PY{o}{.}\PY{n}{mean}\PY{p}{(}\PY{p}{(}\PY{n}{y\PYZus{}valid} \PY{o}{\PYZhy{}} \PY{n}{valid\PYZus{}pred}\PY{p}{)}\PY{o}{*}\PY{o}{*}\PY{l+m+mi}{2}\PY{p}{)}
\end{Verbatim}
\end{tcolorbox}

            \begin{tcolorbox}[breakable, size=fbox, boxrule=.5pt, pad at break*=1mm, opacityfill=0]
\prompt{Out}{outcolor}{5}{\boxspacing}
\begin{Verbatim}[commandchars=\\\{\}]
23.61661706966988
\end{Verbatim}
\end{tcolorbox}
        
    Hence our estimate for the validation MSE of the linear regression fit
is \(23.62\).

We can also estimate the validation error for higher-degree polynomial
regressions. We first provide a function \texttt{evalMSE()} that takes a
model string as well as training and test sets and returns the MSE on
the test set.

    \begin{tcolorbox}[breakable, size=fbox, boxrule=1pt, pad at break*=1mm,colback=cellbackground, colframe=cellborder]
\prompt{In}{incolor}{6}{\boxspacing}
\begin{Verbatim}[commandchars=\\\{\}]
\PY{k}{def}\PY{+w}{ }\PY{n+nf}{evalMSE}\PY{p}{(}\PY{n}{terms}\PY{p}{,}
            \PY{n}{response}\PY{p}{,}
            \PY{n}{train}\PY{p}{,}
            \PY{n}{test}\PY{p}{)}\PY{p}{:}

   \PY{n}{mm} \PY{o}{=} \PY{n}{MS}\PY{p}{(}\PY{n}{terms}\PY{p}{)}
   \PY{n}{X\PYZus{}train} \PY{o}{=} \PY{n}{mm}\PY{o}{.}\PY{n}{fit\PYZus{}transform}\PY{p}{(}\PY{n}{train}\PY{p}{)}
   \PY{n}{y\PYZus{}train} \PY{o}{=} \PY{n}{train}\PY{p}{[}\PY{n}{response}\PY{p}{]}

   \PY{n}{X\PYZus{}test} \PY{o}{=} \PY{n}{mm}\PY{o}{.}\PY{n}{transform}\PY{p}{(}\PY{n}{test}\PY{p}{)}
   \PY{n}{y\PYZus{}test} \PY{o}{=} \PY{n}{test}\PY{p}{[}\PY{n}{response}\PY{p}{]}

   \PY{n}{results} \PY{o}{=} \PY{n}{sm}\PY{o}{.}\PY{n}{OLS}\PY{p}{(}\PY{n}{y\PYZus{}train}\PY{p}{,} \PY{n}{X\PYZus{}train}\PY{p}{)}\PY{o}{.}\PY{n}{fit}\PY{p}{(}\PY{p}{)}
   \PY{n}{test\PYZus{}pred} \PY{o}{=} \PY{n}{results}\PY{o}{.}\PY{n}{predict}\PY{p}{(}\PY{n}{X\PYZus{}test}\PY{p}{)}

   \PY{k}{return} \PY{n}{np}\PY{o}{.}\PY{n}{mean}\PY{p}{(}\PY{p}{(}\PY{n}{y\PYZus{}test} \PY{o}{\PYZhy{}} \PY{n}{test\PYZus{}pred}\PY{p}{)}\PY{o}{*}\PY{o}{*}\PY{l+m+mi}{2}\PY{p}{)}
\end{Verbatim}
\end{tcolorbox}

    Let's use this function to estimate the validation MSE using linear,
quadratic and cubic fits. We use the \texttt{enumerate()} function here,
which gives both the values and indices of objects as one iterates over
a for loop.

    \begin{tcolorbox}[breakable, size=fbox, boxrule=1pt, pad at break*=1mm,colback=cellbackground, colframe=cellborder]
\prompt{In}{incolor}{7}{\boxspacing}
\begin{Verbatim}[commandchars=\\\{\}]
\PY{n}{MSE} \PY{o}{=} \PY{n}{np}\PY{o}{.}\PY{n}{zeros}\PY{p}{(}\PY{l+m+mi}{3}\PY{p}{)}
\PY{k}{for} \PY{n}{idx}\PY{p}{,} \PY{n}{degree} \PY{o+ow}{in} \PY{n+nb}{enumerate}\PY{p}{(}\PY{n+nb}{range}\PY{p}{(}\PY{l+m+mi}{1}\PY{p}{,} \PY{l+m+mi}{4}\PY{p}{)}\PY{p}{)}\PY{p}{:}
    \PY{n}{MSE}\PY{p}{[}\PY{n}{idx}\PY{p}{]} \PY{o}{=} \PY{n}{evalMSE}\PY{p}{(}\PY{p}{[}\PY{n}{poly}\PY{p}{(}\PY{l+s+s1}{\PYZsq{}}\PY{l+s+s1}{horsepower}\PY{l+s+s1}{\PYZsq{}}\PY{p}{,} \PY{n}{degree}\PY{p}{)}\PY{p}{]}\PY{p}{,}
                       \PY{l+s+s1}{\PYZsq{}}\PY{l+s+s1}{mpg}\PY{l+s+s1}{\PYZsq{}}\PY{p}{,}
                       \PY{n}{Auto\PYZus{}train}\PY{p}{,}
                       \PY{n}{Auto\PYZus{}valid}\PY{p}{)}
\PY{n}{MSE}
\end{Verbatim}
\end{tcolorbox}

            \begin{tcolorbox}[breakable, size=fbox, boxrule=.5pt, pad at break*=1mm, opacityfill=0]
\prompt{Out}{outcolor}{7}{\boxspacing}
\begin{Verbatim}[commandchars=\\\{\}]
array([23.61661707, 18.76303135, 18.79694163])
\end{Verbatim}
\end{tcolorbox}
        
    These error rates are \(23.62, 18.76\), and \(18.80\), respectively. If
we choose a different training/validation split instead, then we can
expect somewhat different errors on the validation set.

    \begin{tcolorbox}[breakable, size=fbox, boxrule=1pt, pad at break*=1mm,colback=cellbackground, colframe=cellborder]
\prompt{In}{incolor}{8}{\boxspacing}
\begin{Verbatim}[commandchars=\\\{\}]
\PY{n}{Auto\PYZus{}train}\PY{p}{,} \PY{n}{Auto\PYZus{}valid} \PY{o}{=} \PY{n}{train\PYZus{}test\PYZus{}split}\PY{p}{(}\PY{n}{Auto}\PY{p}{,}
                                          \PY{n}{test\PYZus{}size}\PY{o}{=}\PY{l+m+mi}{196}\PY{p}{,}
                                          \PY{n}{random\PYZus{}state}\PY{o}{=}\PY{l+m+mi}{3}\PY{p}{)}
\PY{n}{MSE} \PY{o}{=} \PY{n}{np}\PY{o}{.}\PY{n}{zeros}\PY{p}{(}\PY{l+m+mi}{3}\PY{p}{)}
\PY{k}{for} \PY{n}{idx}\PY{p}{,} \PY{n}{degree} \PY{o+ow}{in} \PY{n+nb}{enumerate}\PY{p}{(}\PY{n+nb}{range}\PY{p}{(}\PY{l+m+mi}{1}\PY{p}{,} \PY{l+m+mi}{4}\PY{p}{)}\PY{p}{)}\PY{p}{:}
    \PY{n}{MSE}\PY{p}{[}\PY{n}{idx}\PY{p}{]} \PY{o}{=} \PY{n}{evalMSE}\PY{p}{(}\PY{p}{[}\PY{n}{poly}\PY{p}{(}\PY{l+s+s1}{\PYZsq{}}\PY{l+s+s1}{horsepower}\PY{l+s+s1}{\PYZsq{}}\PY{p}{,} \PY{n}{degree}\PY{p}{)}\PY{p}{]}\PY{p}{,}
                       \PY{l+s+s1}{\PYZsq{}}\PY{l+s+s1}{mpg}\PY{l+s+s1}{\PYZsq{}}\PY{p}{,}
                       \PY{n}{Auto\PYZus{}train}\PY{p}{,}
                       \PY{n}{Auto\PYZus{}valid}\PY{p}{)}
\PY{n}{MSE}
\end{Verbatim}
\end{tcolorbox}

            \begin{tcolorbox}[breakable, size=fbox, boxrule=.5pt, pad at break*=1mm, opacityfill=0]
\prompt{Out}{outcolor}{8}{\boxspacing}
\begin{Verbatim}[commandchars=\\\{\}]
array([20.75540796, 16.94510676, 16.97437833])
\end{Verbatim}
\end{tcolorbox}
        
    Using this split of the observations into a training set and a
validation set, we find that the validation set error rates for the
models with linear, quadratic, and cubic terms are \(20.76\), \(16.95\),
and \(16.97\), respectively.

These results are consistent with our previous findings: a model that
predicts \texttt{mpg} using a quadratic function of \texttt{horsepower}
performs better than a model that involves only a linear function of
\texttt{horsepower}, and there is no evidence of an improvement in using
a cubic function of \texttt{horsepower}.

    \hypertarget{cross-validation}{%
\subsection{Cross-Validation}\label{cross-validation}}

In theory, the cross-validation estimate can be computed for any
generalized linear model. \{\} In practice, however, the simplest way to
cross-validate in Python is to use \texttt{sklearn}, which has a
different interface or API than \texttt{statsmodels}, the code we have
been using to fit GLMs.

This is a problem which often confronts data scientists: ``I have a
function to do task \(A\), and need to feed it into something that
performs task \(B\), so that I can compute \(B(A(D))\), where \(D\) is
my data.'' When \(A\) and \(B\) don't naturally speak to each other,
this requires the use of a \emph{wrapper}. In the \texttt{ISLP} package,
we provide a wrapper, \texttt{sklearn\_sm()}, that enables us to easily
use the cross-validation tools of \texttt{sklearn} with models fit by
\texttt{statsmodels}.

The class \texttt{sklearn\_sm()} has as its first argument a model from
\texttt{statsmodels}. It can take two additional optional arguments:
\texttt{model\_str} which can be used to specify a formula, and
\texttt{model\_args} which should be a dictionary of additional
arguments used when fitting the model. For example, to fit a logistic
regression model we have to specify a \texttt{family} argument. This is
passed as
\texttt{model\_args=\{\textquotesingle{}family\textquotesingle{}:sm.families.Binomial()\}}.

Here is our wrapper in action:

    \begin{tcolorbox}[breakable, size=fbox, boxrule=1pt, pad at break*=1mm,colback=cellbackground, colframe=cellborder]
\prompt{In}{incolor}{9}{\boxspacing}
\begin{Verbatim}[commandchars=\\\{\}]
\PY{n}{hp\PYZus{}model} \PY{o}{=} \PY{n}{sklearn\PYZus{}sm}\PY{p}{(}\PY{n}{sm}\PY{o}{.}\PY{n}{OLS}\PY{p}{,}
                      \PY{n}{MS}\PY{p}{(}\PY{p}{[}\PY{l+s+s1}{\PYZsq{}}\PY{l+s+s1}{horsepower}\PY{l+s+s1}{\PYZsq{}}\PY{p}{]}\PY{p}{)}\PY{p}{)}
\PY{n}{X}\PY{p}{,} \PY{n}{Y} \PY{o}{=} \PY{n}{Auto}\PY{o}{.}\PY{n}{drop}\PY{p}{(}\PY{n}{columns}\PY{o}{=}\PY{p}{[}\PY{l+s+s1}{\PYZsq{}}\PY{l+s+s1}{mpg}\PY{l+s+s1}{\PYZsq{}}\PY{p}{]}\PY{p}{)}\PY{p}{,} \PY{n}{Auto}\PY{p}{[}\PY{l+s+s1}{\PYZsq{}}\PY{l+s+s1}{mpg}\PY{l+s+s1}{\PYZsq{}}\PY{p}{]}
\PY{n}{cv\PYZus{}results} \PY{o}{=} \PY{n}{cross\PYZus{}validate}\PY{p}{(}\PY{n}{hp\PYZus{}model}\PY{p}{,}
                            \PY{n}{X}\PY{p}{,}
                            \PY{n}{Y}\PY{p}{,}
                            \PY{n}{cv}\PY{o}{=}\PY{n}{Auto}\PY{o}{.}\PY{n}{shape}\PY{p}{[}\PY{l+m+mi}{0}\PY{p}{]}\PY{p}{)}
\PY{n}{cv\PYZus{}err} \PY{o}{=} \PY{n}{np}\PY{o}{.}\PY{n}{mean}\PY{p}{(}\PY{n}{cv\PYZus{}results}\PY{p}{[}\PY{l+s+s1}{\PYZsq{}}\PY{l+s+s1}{test\PYZus{}score}\PY{l+s+s1}{\PYZsq{}}\PY{p}{]}\PY{p}{)}
\PY{n}{cv\PYZus{}err}
\end{Verbatim}
\end{tcolorbox}

            \begin{tcolorbox}[breakable, size=fbox, boxrule=.5pt, pad at break*=1mm, opacityfill=0]
\prompt{Out}{outcolor}{9}{\boxspacing}
\begin{Verbatim}[commandchars=\\\{\}]
24.23151351792924
\end{Verbatim}
\end{tcolorbox}
        
    The arguments to \texttt{cross\_validate()} are as follows: an object
with the appropriate \texttt{fit()}, \texttt{predict()}, and
\texttt{score()} methods, an array of features \texttt{X} and a response
\texttt{Y}. We also included an additional argument \texttt{cv} to
\texttt{cross\_validate()}; specifying an integer \(k\) results in
\(k\)-fold cross-validation. We have provided a value corresponding to
the total number of observations, which results in leave-one-out
cross-validation (LOOCV). The \texttt{cross\_validate()} function
produces a dictionary with several components; we simply want the
cross-validated test score here (MSE), which is estimated to be 24.23.

    We can repeat this procedure for increasingly complex polynomial fits.
To automate the process, we again use a for loop which iteratively fits
polynomial regressions of degree 1 to 5, computes the associated
cross-validation error, and stores it in the \(i\)th element of the
vector \texttt{cv\_error}. The variable \texttt{d} in the for loop
corresponds to the degree of the polynomial. We begin by initializing
the vector. This command may take a couple of seconds to run.

    \begin{tcolorbox}[breakable, size=fbox, boxrule=1pt, pad at break*=1mm,colback=cellbackground, colframe=cellborder]
\prompt{In}{incolor}{10}{\boxspacing}
\begin{Verbatim}[commandchars=\\\{\}]
\PY{n}{cv\PYZus{}error} \PY{o}{=} \PY{n}{np}\PY{o}{.}\PY{n}{zeros}\PY{p}{(}\PY{l+m+mi}{5}\PY{p}{)}
\PY{n}{H} \PY{o}{=} \PY{n}{np}\PY{o}{.}\PY{n}{array}\PY{p}{(}\PY{n}{Auto}\PY{p}{[}\PY{l+s+s1}{\PYZsq{}}\PY{l+s+s1}{horsepower}\PY{l+s+s1}{\PYZsq{}}\PY{p}{]}\PY{p}{)}
\PY{n}{M} \PY{o}{=} \PY{n}{sklearn\PYZus{}sm}\PY{p}{(}\PY{n}{sm}\PY{o}{.}\PY{n}{OLS}\PY{p}{)}
\PY{k}{for} \PY{n}{i}\PY{p}{,} \PY{n}{d} \PY{o+ow}{in} \PY{n+nb}{enumerate}\PY{p}{(}\PY{n+nb}{range}\PY{p}{(}\PY{l+m+mi}{1}\PY{p}{,}\PY{l+m+mi}{6}\PY{p}{)}\PY{p}{)}\PY{p}{:}
    \PY{n}{X} \PY{o}{=} \PY{n}{np}\PY{o}{.}\PY{n}{power}\PY{o}{.}\PY{n}{outer}\PY{p}{(}\PY{n}{H}\PY{p}{,} \PY{n}{np}\PY{o}{.}\PY{n}{arange}\PY{p}{(}\PY{n}{d}\PY{o}{+}\PY{l+m+mi}{1}\PY{p}{)}\PY{p}{)}
    \PY{n}{M\PYZus{}CV} \PY{o}{=} \PY{n}{cross\PYZus{}validate}\PY{p}{(}\PY{n}{M}\PY{p}{,}
                          \PY{n}{X}\PY{p}{,}
                          \PY{n}{Y}\PY{p}{,}
                          \PY{n}{cv}\PY{o}{=}\PY{n}{Auto}\PY{o}{.}\PY{n}{shape}\PY{p}{[}\PY{l+m+mi}{0}\PY{p}{]}\PY{p}{)}
    \PY{n}{cv\PYZus{}error}\PY{p}{[}\PY{n}{i}\PY{p}{]} \PY{o}{=} \PY{n}{np}\PY{o}{.}\PY{n}{mean}\PY{p}{(}\PY{n}{M\PYZus{}CV}\PY{p}{[}\PY{l+s+s1}{\PYZsq{}}\PY{l+s+s1}{test\PYZus{}score}\PY{l+s+s1}{\PYZsq{}}\PY{p}{]}\PY{p}{)}
\PY{n}{cv\PYZus{}error}
\end{Verbatim}
\end{tcolorbox}

            \begin{tcolorbox}[breakable, size=fbox, boxrule=.5pt, pad at break*=1mm, opacityfill=0]
\prompt{Out}{outcolor}{10}{\boxspacing}
\begin{Verbatim}[commandchars=\\\{\}]
array([24.23151352, 19.24821312, 19.33498406, 19.42443033, 19.03323827])
\end{Verbatim}
\end{tcolorbox}
        
    As in Figure 5.4, we see a sharp drop in the estimated test MSE between
the linear and quadratic fits, but then no clear improvement from using
higher-degree polynomials.

Above we introduced the \texttt{outer()} method of the
\texttt{np.power()} function. The \texttt{outer()} method is applied to
an operation that has two arguments, such as \texttt{add()},
\texttt{min()}, or \texttt{power()}. It has two arrays as arguments, and
then forms a larger array where the operation is applied to each pair of
elements of the two arrays.

    \begin{tcolorbox}[breakable, size=fbox, boxrule=1pt, pad at break*=1mm,colback=cellbackground, colframe=cellborder]
\prompt{In}{incolor}{11}{\boxspacing}
\begin{Verbatim}[commandchars=\\\{\}]
\PY{n}{A} \PY{o}{=} \PY{n}{np}\PY{o}{.}\PY{n}{array}\PY{p}{(}\PY{p}{[}\PY{l+m+mi}{3}\PY{p}{,} \PY{l+m+mi}{5}\PY{p}{,} \PY{l+m+mi}{9}\PY{p}{]}\PY{p}{)}
\PY{n}{B} \PY{o}{=} \PY{n}{np}\PY{o}{.}\PY{n}{array}\PY{p}{(}\PY{p}{[}\PY{l+m+mi}{2}\PY{p}{,} \PY{l+m+mi}{4}\PY{p}{]}\PY{p}{)}
\PY{n}{np}\PY{o}{.}\PY{n}{add}\PY{o}{.}\PY{n}{outer}\PY{p}{(}\PY{n}{A}\PY{p}{,} \PY{n}{B}\PY{p}{)}
\end{Verbatim}
\end{tcolorbox}

            \begin{tcolorbox}[breakable, size=fbox, boxrule=.5pt, pad at break*=1mm, opacityfill=0]
\prompt{Out}{outcolor}{11}{\boxspacing}
\begin{Verbatim}[commandchars=\\\{\}]
array([[ 5,  7],
       [ 7,  9],
       [11, 13]])
\end{Verbatim}
\end{tcolorbox}
        
    In the CV example above, we used \(k=n\), but of course we can also use
\(k<n\). The code is very similar to the above (and is significantly
faster). Here we use \texttt{KFold()} to partition the data into
\(k=10\) random groups. We use \texttt{random\_state} to set a random
seed and initialize a vector \texttt{cv\_error} in which we will store
the CV errors corresponding to the polynomial fits of degrees one to
five.

    \begin{tcolorbox}[breakable, size=fbox, boxrule=1pt, pad at break*=1mm,colback=cellbackground, colframe=cellborder]
\prompt{In}{incolor}{12}{\boxspacing}
\begin{Verbatim}[commandchars=\\\{\}]
\PY{n}{cv\PYZus{}error} \PY{o}{=} \PY{n}{np}\PY{o}{.}\PY{n}{zeros}\PY{p}{(}\PY{l+m+mi}{5}\PY{p}{)}
\PY{n}{cv} \PY{o}{=} \PY{n}{KFold}\PY{p}{(}\PY{n}{n\PYZus{}splits}\PY{o}{=}\PY{l+m+mi}{10}\PY{p}{,}
           \PY{n}{shuffle}\PY{o}{=}\PY{k+kc}{True}\PY{p}{,}
           \PY{n}{random\PYZus{}state}\PY{o}{=}\PY{l+m+mi}{0}\PY{p}{)} \PY{c+c1}{\PYZsh{} use same splits for each degree}
\PY{k}{for} \PY{n}{i}\PY{p}{,} \PY{n}{d} \PY{o+ow}{in} \PY{n+nb}{enumerate}\PY{p}{(}\PY{n+nb}{range}\PY{p}{(}\PY{l+m+mi}{1}\PY{p}{,}\PY{l+m+mi}{6}\PY{p}{)}\PY{p}{)}\PY{p}{:}
    \PY{n}{X} \PY{o}{=} \PY{n}{np}\PY{o}{.}\PY{n}{power}\PY{o}{.}\PY{n}{outer}\PY{p}{(}\PY{n}{H}\PY{p}{,} \PY{n}{np}\PY{o}{.}\PY{n}{arange}\PY{p}{(}\PY{n}{d}\PY{o}{+}\PY{l+m+mi}{1}\PY{p}{)}\PY{p}{)}
    \PY{n}{M\PYZus{}CV} \PY{o}{=} \PY{n}{cross\PYZus{}validate}\PY{p}{(}\PY{n}{M}\PY{p}{,}
                          \PY{n}{X}\PY{p}{,}
                          \PY{n}{Y}\PY{p}{,}
                          \PY{n}{cv}\PY{o}{=}\PY{n}{cv}\PY{p}{)}
    \PY{n}{cv\PYZus{}error}\PY{p}{[}\PY{n}{i}\PY{p}{]} \PY{o}{=} \PY{n}{np}\PY{o}{.}\PY{n}{mean}\PY{p}{(}\PY{n}{M\PYZus{}CV}\PY{p}{[}\PY{l+s+s1}{\PYZsq{}}\PY{l+s+s1}{test\PYZus{}score}\PY{l+s+s1}{\PYZsq{}}\PY{p}{]}\PY{p}{)}
\PY{n}{cv\PYZus{}error}
\end{Verbatim}
\end{tcolorbox}

            \begin{tcolorbox}[breakable, size=fbox, boxrule=.5pt, pad at break*=1mm, opacityfill=0]
\prompt{Out}{outcolor}{12}{\boxspacing}
\begin{Verbatim}[commandchars=\\\{\}]
array([24.20766449, 19.18533142, 19.27626666, 19.47848403, 19.13720581])
\end{Verbatim}
\end{tcolorbox}
        
    Notice that the computation time is much shorter than that of LOOCV. (In
principle, the computation time for LOOCV for a least squares linear
model should be faster than for \(k\)-fold CV, due to the availability
of the formula\textasciitilde(5.2) for LOOCV; however, the generic
\texttt{cross\_validate()} function does not make use of this formula.)
We still see little evidence that using cubic or higher-degree
polynomial terms leads to a lower test error than simply using a
quadratic fit.

    The \texttt{cross\_validate()} function is flexible and can take
different splitting mechanisms as an argument. For instance, one can use
the \texttt{ShuffleSplit()} function to implement the validation set
approach just as easily as \(k\)-fold cross-validation.

    \begin{tcolorbox}[breakable, size=fbox, boxrule=1pt, pad at break*=1mm,colback=cellbackground, colframe=cellborder]
\prompt{In}{incolor}{13}{\boxspacing}
\begin{Verbatim}[commandchars=\\\{\}]
\PY{n}{validation} \PY{o}{=} \PY{n}{ShuffleSplit}\PY{p}{(}\PY{n}{n\PYZus{}splits}\PY{o}{=}\PY{l+m+mi}{1}\PY{p}{,}
                          \PY{n}{test\PYZus{}size}\PY{o}{=}\PY{l+m+mi}{196}\PY{p}{,}
                          \PY{n}{random\PYZus{}state}\PY{o}{=}\PY{l+m+mi}{0}\PY{p}{)}
\PY{n}{results} \PY{o}{=} \PY{n}{cross\PYZus{}validate}\PY{p}{(}\PY{n}{hp\PYZus{}model}\PY{p}{,}
                         \PY{n}{Auto}\PY{o}{.}\PY{n}{drop}\PY{p}{(}\PY{p}{[}\PY{l+s+s1}{\PYZsq{}}\PY{l+s+s1}{mpg}\PY{l+s+s1}{\PYZsq{}}\PY{p}{]}\PY{p}{,} \PY{n}{axis}\PY{o}{=}\PY{l+m+mi}{1}\PY{p}{)}\PY{p}{,}
                         \PY{n}{Auto}\PY{p}{[}\PY{l+s+s1}{\PYZsq{}}\PY{l+s+s1}{mpg}\PY{l+s+s1}{\PYZsq{}}\PY{p}{]}\PY{p}{,}
                         \PY{n}{cv}\PY{o}{=}\PY{n}{validation}\PY{p}{)}\PY{p}{;}
\PY{n}{results}\PY{p}{[}\PY{l+s+s1}{\PYZsq{}}\PY{l+s+s1}{test\PYZus{}score}\PY{l+s+s1}{\PYZsq{}}\PY{p}{]}
\end{Verbatim}
\end{tcolorbox}

            \begin{tcolorbox}[breakable, size=fbox, boxrule=.5pt, pad at break*=1mm, opacityfill=0]
\prompt{Out}{outcolor}{13}{\boxspacing}
\begin{Verbatim}[commandchars=\\\{\}]
array([23.61661707])
\end{Verbatim}
\end{tcolorbox}
        
    One can estimate the variability in the test error by running the
following:

    \begin{tcolorbox}[breakable, size=fbox, boxrule=1pt, pad at break*=1mm,colback=cellbackground, colframe=cellborder]
\prompt{In}{incolor}{14}{\boxspacing}
\begin{Verbatim}[commandchars=\\\{\}]
\PY{n}{validation} \PY{o}{=} \PY{n}{ShuffleSplit}\PY{p}{(}\PY{n}{n\PYZus{}splits}\PY{o}{=}\PY{l+m+mi}{10}\PY{p}{,}
                          \PY{n}{test\PYZus{}size}\PY{o}{=}\PY{l+m+mi}{196}\PY{p}{,}
                          \PY{n}{random\PYZus{}state}\PY{o}{=}\PY{l+m+mi}{0}\PY{p}{)}
\PY{n}{results} \PY{o}{=} \PY{n}{cross\PYZus{}validate}\PY{p}{(}\PY{n}{hp\PYZus{}model}\PY{p}{,}
                         \PY{n}{Auto}\PY{o}{.}\PY{n}{drop}\PY{p}{(}\PY{p}{[}\PY{l+s+s1}{\PYZsq{}}\PY{l+s+s1}{mpg}\PY{l+s+s1}{\PYZsq{}}\PY{p}{]}\PY{p}{,} \PY{n}{axis}\PY{o}{=}\PY{l+m+mi}{1}\PY{p}{)}\PY{p}{,}
                         \PY{n}{Auto}\PY{p}{[}\PY{l+s+s1}{\PYZsq{}}\PY{l+s+s1}{mpg}\PY{l+s+s1}{\PYZsq{}}\PY{p}{]}\PY{p}{,}
                         \PY{n}{cv}\PY{o}{=}\PY{n}{validation}\PY{p}{)}
\PY{n}{results}\PY{p}{[}\PY{l+s+s1}{\PYZsq{}}\PY{l+s+s1}{test\PYZus{}score}\PY{l+s+s1}{\PYZsq{}}\PY{p}{]}\PY{o}{.}\PY{n}{mean}\PY{p}{(}\PY{p}{)}\PY{p}{,} \PY{n}{results}\PY{p}{[}\PY{l+s+s1}{\PYZsq{}}\PY{l+s+s1}{test\PYZus{}score}\PY{l+s+s1}{\PYZsq{}}\PY{p}{]}\PY{o}{.}\PY{n}{std}\PY{p}{(}\PY{p}{)}
\end{Verbatim}
\end{tcolorbox}

            \begin{tcolorbox}[breakable, size=fbox, boxrule=.5pt, pad at break*=1mm, opacityfill=0]
\prompt{Out}{outcolor}{14}{\boxspacing}
\begin{Verbatim}[commandchars=\\\{\}]
(23.802232661034168, 1.421845094109185)
\end{Verbatim}
\end{tcolorbox}
        
    Note that this standard deviation is not a valid estimate of the
sampling variability of the mean test score or the individual scores,
since the randomly-selected training samples overlap and hence introduce
correlations. But it does give an idea of the Monte Carlo variation
incurred by picking different random folds.

\hypertarget{the-bootstrap}{%
\subsection{The Bootstrap}\label{the-bootstrap}}

We illustrate the use of the bootstrap in the simple example \{of
Section 5.2,\} as well as on an example involving estimating the
accuracy of the linear regression model on the \texttt{Auto} data set.
\#\#\# Estimating the Accuracy of a Statistic of Interest One of the
great advantages of the bootstrap approach is that it can be applied in
almost all situations. No complicated mathematical calculations are
required. While there are several implementations of the bootstrap in
Python, its use for estimating standard error is simple enough that we
write our own function below for the case when our data is stored in a
dataframe.

To illustrate the bootstrap, we start with a simple example. The
\texttt{Portfolio} data set in the \texttt{ISLP} package is described in
Section 5.2. The goal is to estimate the sampling variance of the
parameter \(\alpha\) given in formula\textasciitilde(5.7). We will
create a function \texttt{alpha\_func()}, which takes as input a
dataframe \texttt{D} assumed to have columns \texttt{X} and \texttt{Y},
as well as a vector \texttt{idx} indicating which observations should be
used to estimate \(\alpha\). The function then outputs the estimate for
\(\alpha\) based on the selected observations.

    \begin{tcolorbox}[breakable, size=fbox, boxrule=1pt, pad at break*=1mm,colback=cellbackground, colframe=cellborder]
\prompt{In}{incolor}{15}{\boxspacing}
\begin{Verbatim}[commandchars=\\\{\}]
\PY{n}{Portfolio} \PY{o}{=} \PY{n}{load\PYZus{}data}\PY{p}{(}\PY{l+s+s1}{\PYZsq{}}\PY{l+s+s1}{Portfolio}\PY{l+s+s1}{\PYZsq{}}\PY{p}{)}
\PY{k}{def}\PY{+w}{ }\PY{n+nf}{alpha\PYZus{}func}\PY{p}{(}\PY{n}{D}\PY{p}{,} \PY{n}{idx}\PY{p}{)}\PY{p}{:}
   \PY{n}{cov\PYZus{}} \PY{o}{=} \PY{n}{np}\PY{o}{.}\PY{n}{cov}\PY{p}{(}\PY{n}{D}\PY{p}{[}\PY{p}{[}\PY{l+s+s1}{\PYZsq{}}\PY{l+s+s1}{X}\PY{l+s+s1}{\PYZsq{}}\PY{p}{,}\PY{l+s+s1}{\PYZsq{}}\PY{l+s+s1}{Y}\PY{l+s+s1}{\PYZsq{}}\PY{p}{]}\PY{p}{]}\PY{o}{.}\PY{n}{loc}\PY{p}{[}\PY{n}{idx}\PY{p}{]}\PY{p}{,} \PY{n}{rowvar}\PY{o}{=}\PY{k+kc}{False}\PY{p}{)}
   \PY{k}{return} \PY{p}{(}\PY{p}{(}\PY{n}{cov\PYZus{}}\PY{p}{[}\PY{l+m+mi}{1}\PY{p}{,}\PY{l+m+mi}{1}\PY{p}{]} \PY{o}{\PYZhy{}} \PY{n}{cov\PYZus{}}\PY{p}{[}\PY{l+m+mi}{0}\PY{p}{,}\PY{l+m+mi}{1}\PY{p}{]}\PY{p}{)} \PY{o}{/}
           \PY{p}{(}\PY{n}{cov\PYZus{}}\PY{p}{[}\PY{l+m+mi}{0}\PY{p}{,}\PY{l+m+mi}{0}\PY{p}{]}\PY{o}{+}\PY{n}{cov\PYZus{}}\PY{p}{[}\PY{l+m+mi}{1}\PY{p}{,}\PY{l+m+mi}{1}\PY{p}{]}\PY{o}{\PYZhy{}}\PY{l+m+mi}{2}\PY{o}{*}\PY{n}{cov\PYZus{}}\PY{p}{[}\PY{l+m+mi}{0}\PY{p}{,}\PY{l+m+mi}{1}\PY{p}{]}\PY{p}{)}\PY{p}{)}
\end{Verbatim}
\end{tcolorbox}

    This function returns an estimate for \(\alpha\) based on applying the
minimum variance formula (5.7) to the observations indexed by the
argument \texttt{idx}. For instance, the following command estimates
\(\alpha\) using all 100 observations.

    \begin{tcolorbox}[breakable, size=fbox, boxrule=1pt, pad at break*=1mm,colback=cellbackground, colframe=cellborder]
\prompt{In}{incolor}{16}{\boxspacing}
\begin{Verbatim}[commandchars=\\\{\}]
\PY{n}{alpha\PYZus{}func}\PY{p}{(}\PY{n}{Portfolio}\PY{p}{,} \PY{n+nb}{range}\PY{p}{(}\PY{l+m+mi}{100}\PY{p}{)}\PY{p}{)}
\end{Verbatim}
\end{tcolorbox}

            \begin{tcolorbox}[breakable, size=fbox, boxrule=.5pt, pad at break*=1mm, opacityfill=0]
\prompt{Out}{outcolor}{16}{\boxspacing}
\begin{Verbatim}[commandchars=\\\{\}]
0.57583207459283
\end{Verbatim}
\end{tcolorbox}
        
    Next we randomly select 100 observations from \texttt{range(100)}, with
replacement. This is equivalent to constructing a new bootstrap data set
and recomputing \(\hat{\alpha}\) based on the new data set.

    \begin{tcolorbox}[breakable, size=fbox, boxrule=1pt, pad at break*=1mm,colback=cellbackground, colframe=cellborder]
\prompt{In}{incolor}{17}{\boxspacing}
\begin{Verbatim}[commandchars=\\\{\}]
\PY{n}{rng} \PY{o}{=} \PY{n}{np}\PY{o}{.}\PY{n}{random}\PY{o}{.}\PY{n}{default\PYZus{}rng}\PY{p}{(}\PY{l+m+mi}{0}\PY{p}{)}
\PY{n}{alpha\PYZus{}func}\PY{p}{(}\PY{n}{Portfolio}\PY{p}{,}
           \PY{n}{rng}\PY{o}{.}\PY{n}{choice}\PY{p}{(}\PY{l+m+mi}{100}\PY{p}{,}
                      \PY{l+m+mi}{100}\PY{p}{,}
                      \PY{n}{replace}\PY{o}{=}\PY{k+kc}{True}\PY{p}{)}\PY{p}{)}
\end{Verbatim}
\end{tcolorbox}

            \begin{tcolorbox}[breakable, size=fbox, boxrule=.5pt, pad at break*=1mm, opacityfill=0]
\prompt{Out}{outcolor}{17}{\boxspacing}
\begin{Verbatim}[commandchars=\\\{\}]
0.6074452469619004
\end{Verbatim}
\end{tcolorbox}
        
    This process can be generalized to create a simple function
\texttt{boot\_SE()} for computing the bootstrap standard error for
arbitrary functions that take only a data frame as an argument.

    \begin{tcolorbox}[breakable, size=fbox, boxrule=1pt, pad at break*=1mm,colback=cellbackground, colframe=cellborder]
\prompt{In}{incolor}{18}{\boxspacing}
\begin{Verbatim}[commandchars=\\\{\}]
\PY{k}{def}\PY{+w}{ }\PY{n+nf}{boot\PYZus{}SE}\PY{p}{(}\PY{n}{func}\PY{p}{,}
            \PY{n}{D}\PY{p}{,}
            \PY{n}{n}\PY{o}{=}\PY{k+kc}{None}\PY{p}{,}
            \PY{n}{B}\PY{o}{=}\PY{l+m+mi}{1000}\PY{p}{,}
            \PY{n}{seed}\PY{o}{=}\PY{l+m+mi}{0}\PY{p}{)}\PY{p}{:}
    \PY{n}{rng} \PY{o}{=} \PY{n}{np}\PY{o}{.}\PY{n}{random}\PY{o}{.}\PY{n}{default\PYZus{}rng}\PY{p}{(}\PY{n}{seed}\PY{p}{)}
    \PY{n}{first\PYZus{}}\PY{p}{,} \PY{n}{second\PYZus{}} \PY{o}{=} \PY{l+m+mi}{0}\PY{p}{,} \PY{l+m+mi}{0}
    \PY{n}{n} \PY{o}{=} \PY{n}{n} \PY{o+ow}{or} \PY{n}{D}\PY{o}{.}\PY{n}{shape}\PY{p}{[}\PY{l+m+mi}{0}\PY{p}{]}
    \PY{k}{for} \PY{n}{\PYZus{}} \PY{o+ow}{in} \PY{n+nb}{range}\PY{p}{(}\PY{n}{B}\PY{p}{)}\PY{p}{:}
        \PY{n}{idx} \PY{o}{=} \PY{n}{rng}\PY{o}{.}\PY{n}{choice}\PY{p}{(}\PY{n}{D}\PY{o}{.}\PY{n}{index}\PY{p}{,}
                         \PY{n}{n}\PY{p}{,}
                         \PY{n}{replace}\PY{o}{=}\PY{k+kc}{True}\PY{p}{)}
        \PY{n}{value} \PY{o}{=} \PY{n}{func}\PY{p}{(}\PY{n}{D}\PY{p}{,} \PY{n}{idx}\PY{p}{)}
        \PY{n}{first\PYZus{}} \PY{o}{+}\PY{o}{=} \PY{n}{value}
        \PY{n}{second\PYZus{}} \PY{o}{+}\PY{o}{=} \PY{n}{value}\PY{o}{*}\PY{o}{*}\PY{l+m+mi}{2}
    \PY{k}{return} \PY{n}{np}\PY{o}{.}\PY{n}{sqrt}\PY{p}{(}\PY{n}{second\PYZus{}} \PY{o}{/} \PY{n}{B} \PY{o}{\PYZhy{}} \PY{p}{(}\PY{n}{first\PYZus{}} \PY{o}{/} \PY{n}{B}\PY{p}{)}\PY{o}{*}\PY{o}{*}\PY{l+m+mi}{2}\PY{p}{)}
\end{Verbatim}
\end{tcolorbox}

    Notice the use of \texttt{\_} as a loop variable in
\texttt{for\ \_\ in\ range(B)}. This is often used if the value of the
counter is unimportant and simply makes sure the loop is executed
\texttt{B} times.

Let's use our function to evaluate the accuracy of our estimate of
\(\alpha\) using \(B=1{,}000\) bootstrap replications.

    \begin{tcolorbox}[breakable, size=fbox, boxrule=1pt, pad at break*=1mm,colback=cellbackground, colframe=cellborder]
\prompt{In}{incolor}{19}{\boxspacing}
\begin{Verbatim}[commandchars=\\\{\}]
\PY{n}{alpha\PYZus{}SE} \PY{o}{=} \PY{n}{boot\PYZus{}SE}\PY{p}{(}\PY{n}{alpha\PYZus{}func}\PY{p}{,}
                   \PY{n}{Portfolio}\PY{p}{,}
                   \PY{n}{B}\PY{o}{=}\PY{l+m+mi}{1000}\PY{p}{,}
                   \PY{n}{seed}\PY{o}{=}\PY{l+m+mi}{0}\PY{p}{)}
\PY{n}{alpha\PYZus{}SE}
\end{Verbatim}
\end{tcolorbox}

            \begin{tcolorbox}[breakable, size=fbox, boxrule=.5pt, pad at break*=1mm, opacityfill=0]
\prompt{Out}{outcolor}{19}{\boxspacing}
\begin{Verbatim}[commandchars=\\\{\}]
0.09118176521277699
\end{Verbatim}
\end{tcolorbox}
        
    The final output shows that the bootstrap estimate for
\({\rm SE}(\hat{\alpha})\) is \(0.0912\).

\hypertarget{estimating-the-accuracy-of-a-linear-regression-model}{%
\subsubsection{Estimating the Accuracy of a Linear Regression
Model}\label{estimating-the-accuracy-of-a-linear-regression-model}}

The bootstrap approach can be used to assess the variability of the
coefficient estimates and predictions from a statistical learning
method. Here we use the bootstrap approach in order to assess the
variability of the estimates for \(\beta_0\) and \(\beta_1\), the
intercept and slope terms for the linear regression model that uses
\texttt{horsepower} to predict \texttt{mpg} in the \texttt{Auto} data
set. We will compare the estimates obtained using the bootstrap to those
obtained using the formulas for \({\rm SE}(\hat{\beta}_0)\) and
\({\rm SE}(\hat{\beta}_1)\) described in Section 3.1.2.

To use our \texttt{boot\_SE()} function, we must write a function (its
first argument) that takes a data frame \texttt{D} and indices
\texttt{idx} as its only arguments. But here we want to bootstrap a
specific regression model, specified by a model formula and data. We
show how to do this in a few simple steps.

We start by writing a generic function \texttt{boot\_OLS()} for
bootstrapping a regression model that takes a formula to define the
corresponding regression. We use the \texttt{clone()} function to make a
copy of the formula that can be refit to the new dataframe. This means
that any derived features such as those defined by \texttt{poly()}
(which we will see shortly), will be re-fit on the resampled data frame.

    \begin{tcolorbox}[breakable, size=fbox, boxrule=1pt, pad at break*=1mm,colback=cellbackground, colframe=cellborder]
\prompt{In}{incolor}{20}{\boxspacing}
\begin{Verbatim}[commandchars=\\\{\}]
\PY{k}{def}\PY{+w}{ }\PY{n+nf}{boot\PYZus{}OLS}\PY{p}{(}\PY{n}{model\PYZus{}matrix}\PY{p}{,} \PY{n}{response}\PY{p}{,} \PY{n}{D}\PY{p}{,} \PY{n}{idx}\PY{p}{)}\PY{p}{:}
    \PY{n}{D\PYZus{}} \PY{o}{=} \PY{n}{D}\PY{o}{.}\PY{n}{loc}\PY{p}{[}\PY{n}{idx}\PY{p}{]}
    \PY{n}{Y\PYZus{}} \PY{o}{=} \PY{n}{D\PYZus{}}\PY{p}{[}\PY{n}{response}\PY{p}{]}
    \PY{n}{X\PYZus{}} \PY{o}{=} \PY{n}{clone}\PY{p}{(}\PY{n}{model\PYZus{}matrix}\PY{p}{)}\PY{o}{.}\PY{n}{fit\PYZus{}transform}\PY{p}{(}\PY{n}{D\PYZus{}}\PY{p}{)}
    \PY{k}{return} \PY{n}{sm}\PY{o}{.}\PY{n}{OLS}\PY{p}{(}\PY{n}{Y\PYZus{}}\PY{p}{,} \PY{n}{X\PYZus{}}\PY{p}{)}\PY{o}{.}\PY{n}{fit}\PY{p}{(}\PY{p}{)}\PY{o}{.}\PY{n}{params}
\end{Verbatim}
\end{tcolorbox}

    This is not quite what is needed as the first argument to
\texttt{boot\_SE()}. The first two arguments which specify the model
will not change in the bootstrap process, and we would like to
\emph{freeze} them. The function \texttt{partial()} from the
\texttt{functools} module does precisely this: it takes a function as an
argument, and freezes some of its arguments, starting from the left. We
use it to freeze the first two model-formula arguments of
\texttt{boot\_OLS()}.

    \begin{tcolorbox}[breakable, size=fbox, boxrule=1pt, pad at break*=1mm,colback=cellbackground, colframe=cellborder]
\prompt{In}{incolor}{21}{\boxspacing}
\begin{Verbatim}[commandchars=\\\{\}]
\PY{n}{hp\PYZus{}func} \PY{o}{=} \PY{n}{partial}\PY{p}{(}\PY{n}{boot\PYZus{}OLS}\PY{p}{,} \PY{n}{MS}\PY{p}{(}\PY{p}{[}\PY{l+s+s1}{\PYZsq{}}\PY{l+s+s1}{horsepower}\PY{l+s+s1}{\PYZsq{}}\PY{p}{]}\PY{p}{)}\PY{p}{,} \PY{l+s+s1}{\PYZsq{}}\PY{l+s+s1}{mpg}\PY{l+s+s1}{\PYZsq{}}\PY{p}{)}
\end{Verbatim}
\end{tcolorbox}

    Typing \texttt{hp\_func?} will show that it has two arguments \texttt{D}
and \texttt{idx} --- it is a version of \texttt{boot\_OLS()} with the
first two arguments frozen --- and hence is ideal as the first argument
for \texttt{boot\_SE()}.

The \texttt{hp\_func()} function can now be used in order to create
bootstrap estimates for the intercept and slope terms by randomly
sampling from among the observations with replacement. We first
demonstrate its utility on 10 bootstrap samples.

    \begin{tcolorbox}[breakable, size=fbox, boxrule=1pt, pad at break*=1mm,colback=cellbackground, colframe=cellborder]
\prompt{In}{incolor}{22}{\boxspacing}
\begin{Verbatim}[commandchars=\\\{\}]
\PY{n}{rng} \PY{o}{=} \PY{n}{np}\PY{o}{.}\PY{n}{random}\PY{o}{.}\PY{n}{default\PYZus{}rng}\PY{p}{(}\PY{l+m+mi}{0}\PY{p}{)}
\PY{n}{np}\PY{o}{.}\PY{n}{array}\PY{p}{(}\PY{p}{[}\PY{n}{hp\PYZus{}func}\PY{p}{(}\PY{n}{Auto}\PY{p}{,}
          \PY{n}{rng}\PY{o}{.}\PY{n}{choice}\PY{p}{(}\PY{n}{Auto}\PY{o}{.}\PY{n}{index}\PY{p}{,}
                     \PY{l+m+mi}{392}\PY{p}{,}
                     \PY{n}{replace}\PY{o}{=}\PY{k+kc}{True}\PY{p}{)}\PY{p}{)} \PY{k}{for} \PY{n}{\PYZus{}} \PY{o+ow}{in} \PY{n+nb}{range}\PY{p}{(}\PY{l+m+mi}{10}\PY{p}{)}\PY{p}{]}\PY{p}{)}
\end{Verbatim}
\end{tcolorbox}

            \begin{tcolorbox}[breakable, size=fbox, boxrule=.5pt, pad at break*=1mm, opacityfill=0]
\prompt{Out}{outcolor}{22}{\boxspacing}
\begin{Verbatim}[commandchars=\\\{\}]
array([[39.12226577, -0.1555926 ],
       [37.18648613, -0.13915813],
       [37.46989244, -0.14112749],
       [38.56723252, -0.14830116],
       [38.95495707, -0.15315141],
       [39.12563927, -0.15261044],
       [38.45763251, -0.14767251],
       [38.43372587, -0.15019447],
       [37.87581142, -0.1409544 ],
       [37.95949036, -0.1451333 ]])
\end{Verbatim}
\end{tcolorbox}
        
    Next, we use the \texttt{boot\_SE()} \{\} function to compute the
standard errors of 1,000 bootstrap estimates for the intercept and slope
terms.

    \begin{tcolorbox}[breakable, size=fbox, boxrule=1pt, pad at break*=1mm,colback=cellbackground, colframe=cellborder]
\prompt{In}{incolor}{23}{\boxspacing}
\begin{Verbatim}[commandchars=\\\{\}]
\PY{n}{hp\PYZus{}se} \PY{o}{=} \PY{n}{boot\PYZus{}SE}\PY{p}{(}\PY{n}{hp\PYZus{}func}\PY{p}{,}
                \PY{n}{Auto}\PY{p}{,}
                \PY{n}{B}\PY{o}{=}\PY{l+m+mi}{1000}\PY{p}{,}
                \PY{n}{seed}\PY{o}{=}\PY{l+m+mi}{10}\PY{p}{)}
\PY{n}{hp\PYZus{}se}
\end{Verbatim}
\end{tcolorbox}

            \begin{tcolorbox}[breakable, size=fbox, boxrule=.5pt, pad at break*=1mm, opacityfill=0]
\prompt{Out}{outcolor}{23}{\boxspacing}
\begin{Verbatim}[commandchars=\\\{\}]
intercept     0.731176
horsepower    0.006092
dtype: float64
\end{Verbatim}
\end{tcolorbox}
        
    This indicates that the bootstrap estimate for
\({\rm SE}(\hat{\beta}_0)\) is 0.85, and that the bootstrap estimate for
\({\rm SE}(\hat{\beta}_1)\) is 0.0074. As discussed in Section 3.1.2,
standard formulas can be used to compute the standard errors for the
regression coefficients in a linear model. These can be obtained using
the \texttt{summarize()} function from \texttt{ISLP.sm}.

    \begin{tcolorbox}[breakable, size=fbox, boxrule=1pt, pad at break*=1mm,colback=cellbackground, colframe=cellborder]
\prompt{In}{incolor}{24}{\boxspacing}
\begin{Verbatim}[commandchars=\\\{\}]
\PY{n}{hp\PYZus{}model}\PY{o}{.}\PY{n}{fit}\PY{p}{(}\PY{n}{Auto}\PY{p}{,} \PY{n}{Auto}\PY{p}{[}\PY{l+s+s1}{\PYZsq{}}\PY{l+s+s1}{mpg}\PY{l+s+s1}{\PYZsq{}}\PY{p}{]}\PY{p}{)}
\PY{n}{model\PYZus{}se} \PY{o}{=} \PY{n}{summarize}\PY{p}{(}\PY{n}{hp\PYZus{}model}\PY{o}{.}\PY{n}{results\PYZus{}}\PY{p}{)}\PY{p}{[}\PY{l+s+s1}{\PYZsq{}}\PY{l+s+s1}{std err}\PY{l+s+s1}{\PYZsq{}}\PY{p}{]}
\PY{n}{model\PYZus{}se}
\end{Verbatim}
\end{tcolorbox}

            \begin{tcolorbox}[breakable, size=fbox, boxrule=.5pt, pad at break*=1mm, opacityfill=0]
\prompt{Out}{outcolor}{24}{\boxspacing}
\begin{Verbatim}[commandchars=\\\{\}]
intercept     0.717
horsepower    0.006
Name: std err, dtype: float64
\end{Verbatim}
\end{tcolorbox}
        
    The standard error estimates for \(\hat{\beta}_0\) and \(\hat{\beta}_1\)
obtained using the formulas from Section 3.1.2 are 0.717 for the
intercept and 0.006 for the slope. Interestingly, these are somewhat
different from the estimates obtained using the bootstrap. Does this
indicate a problem with the bootstrap? In fact, it suggests the
opposite. Recall that the standard formulas given in \{Equation 3.8 on
page 75\} rely on certain assumptions. For example, they depend on the
unknown parameter \(\sigma^2\), the noise variance. We then estimate
\(\sigma^2\) using the RSS. Now although the formulas for the standard
errors do not rely on the linear model being correct, the estimate for
\(\sigma^2\) does. We see \{in Figure 3.8 on page 99\} that there is a
non-linear relationship in the data, and so the residuals from a linear
fit will be inflated, and so will \(\hat{\sigma}^2\). Secondly, the
standard formulas assume (somewhat unrealistically) that the \(x_i\) are
fixed, and all the variability comes from the variation in the errors
\(\epsilon_i\). The bootstrap approach does not rely on any of these
assumptions, and so it is likely giving a more accurate estimate of the
standard errors of \(\hat{\beta}_0\) and \(\hat{\beta}_1\) than the
results from \texttt{sm.OLS}.

Below we compute the bootstrap standard error estimates and the standard
linear regression estimates that result from fitting the quadratic model
to the data. Since this model provides a good fit to the data (Figure
3.8), there is now a better correspondence between the bootstrap
estimates and the standard estimates of \({\rm SE}(\hat{\beta}_0)\),
\({\rm SE}(\hat{\beta}_1)\) and \({\rm SE}(\hat{\beta}_2)\).

    \begin{tcolorbox}[breakable, size=fbox, boxrule=1pt, pad at break*=1mm,colback=cellbackground, colframe=cellborder]
\prompt{In}{incolor}{25}{\boxspacing}
\begin{Verbatim}[commandchars=\\\{\}]
\PY{n}{quad\PYZus{}model} \PY{o}{=} \PY{n}{MS}\PY{p}{(}\PY{p}{[}\PY{n}{poly}\PY{p}{(}\PY{l+s+s1}{\PYZsq{}}\PY{l+s+s1}{horsepower}\PY{l+s+s1}{\PYZsq{}}\PY{p}{,} \PY{l+m+mi}{2}\PY{p}{,} \PY{n}{raw}\PY{o}{=}\PY{k+kc}{True}\PY{p}{)}\PY{p}{]}\PY{p}{)}
\PY{n}{quad\PYZus{}func} \PY{o}{=} \PY{n}{partial}\PY{p}{(}\PY{n}{boot\PYZus{}OLS}\PY{p}{,}
                    \PY{n}{quad\PYZus{}model}\PY{p}{,}
                    \PY{l+s+s1}{\PYZsq{}}\PY{l+s+s1}{mpg}\PY{l+s+s1}{\PYZsq{}}\PY{p}{)}
\PY{n}{boot\PYZus{}SE}\PY{p}{(}\PY{n}{quad\PYZus{}func}\PY{p}{,} \PY{n}{Auto}\PY{p}{,} \PY{n}{B}\PY{o}{=}\PY{l+m+mi}{1000}\PY{p}{)}
\end{Verbatim}
\end{tcolorbox}

            \begin{tcolorbox}[breakable, size=fbox, boxrule=.5pt, pad at break*=1mm, opacityfill=0]
\prompt{Out}{outcolor}{25}{\boxspacing}
\begin{Verbatim}[commandchars=\\\{\}]
intercept                                  1.538641
poly(horsepower, degree=2, raw=True)[0]    0.024696
poly(horsepower, degree=2, raw=True)[1]    0.000090
dtype: float64
\end{Verbatim}
\end{tcolorbox}
        
    We compare the results to the standard errors computed using
\texttt{sm.OLS()}.

    \begin{tcolorbox}[breakable, size=fbox, boxrule=1pt, pad at break*=1mm,colback=cellbackground, colframe=cellborder]
\prompt{In}{incolor}{26}{\boxspacing}
\begin{Verbatim}[commandchars=\\\{\}]
\PY{n}{M} \PY{o}{=} \PY{n}{sm}\PY{o}{.}\PY{n}{OLS}\PY{p}{(}\PY{n}{Auto}\PY{p}{[}\PY{l+s+s1}{\PYZsq{}}\PY{l+s+s1}{mpg}\PY{l+s+s1}{\PYZsq{}}\PY{p}{]}\PY{p}{,}
           \PY{n}{quad\PYZus{}model}\PY{o}{.}\PY{n}{fit\PYZus{}transform}\PY{p}{(}\PY{n}{Auto}\PY{p}{)}\PY{p}{)}
\PY{n}{summarize}\PY{p}{(}\PY{n}{M}\PY{o}{.}\PY{n}{fit}\PY{p}{(}\PY{p}{)}\PY{p}{)}\PY{p}{[}\PY{l+s+s1}{\PYZsq{}}\PY{l+s+s1}{std err}\PY{l+s+s1}{\PYZsq{}}\PY{p}{]}
\end{Verbatim}
\end{tcolorbox}

            \begin{tcolorbox}[breakable, size=fbox, boxrule=.5pt, pad at break*=1mm, opacityfill=0]
\prompt{Out}{outcolor}{26}{\boxspacing}
\begin{Verbatim}[commandchars=\\\{\}]
intercept                                  1.800
poly(horsepower, degree=2, raw=True)[0]    0.031
poly(horsepower, degree=2, raw=True)[1]    0.000
Name: std err, dtype: float64
\end{Verbatim}
\end{tcolorbox}
        
    


    % Add a bibliography block to the postdoc
    
    
    
\end{document}
